\documentclass[a4paper,12pt]{article}

\usepackage{fontspec}
\usepackage{xunicode}
\usepackage{xltxtra}
\usepackage{graphicx}
\usepackage{float}
\usepackage{hyperref}
\usepackage{geometry}
\usepackage{booktabs}
\usepackage{longtable}

% ตั้งค่าภาษาไทยสำหรับ macOS
\setmainfont{Thonburi}
\XeTeXlinebreaklocale "th"
\XeTeXlinebreakskip = 0pt plus 1pt

\geometry{
 a4paper,
 total={170mm,257mm},
 left=20mm,
 top=20mm,
}

\title{\textbf{ระบบตรวจจับด้วยการเรียนรู้เชิงลึกสำหรับระบุเนื้อปลาแซลมอนและปลาเรนโบว์เทราต์\\และการเกรดความสดเพื่อคุ้มครองผู้บริโภค}}
\author{สรุปและเรียบเรียงจากงานวิจัย PMC12567205}
\date{\today}

\begin{document}

\maketitle

\section{บทคัดย่อ (Abstract)}
การฉ้อโกงอาหารทะเล เช่น การนำปลาเรนโบว์เทราต์ (Rainbow Trout) ราคาถูกมาปลอมแปลงฉลากเป็นปลาแซลมอน (Salmon) เกรดพรีเมียม ก่อให้เกิดความเสี่ยงด้านความปลอดภัยทางอาหารและละเมิดสิทธิผู้บริโภค งานวิจัยนี้ได้พัฒนาระบบตรวจจับที่ทำงานบนสมาร์ทโฟนโดยอาศัยเทคโนโลยีการเรียนรู้เชิงลึก (Deep Learning) เพื่อระบุชนิดของเนื้อปลาและเกรดความสดของปลาแซลมอน

ระบบนี้ใช้การออกแบบแบบสองขั้นตอน (Two-stage design):
\begin{enumerate}
    \item จำแนกชนิดของเนื้อปลา
    \item ให้เกรดความสดของปลาแซลมอนเป็น 3 ระดับ
\end{enumerate}

งานวิจัยใช้สถาปัตยกรรม \textbf{DenseNet121} ที่ได้รับการปรับปรุง (Improved DenseNet121) โดยเพิ่ม Global Average Pooling, Dropout layers และปรับแต่ง Output layer เพื่อเพิ่มความแม่นยำและลดการเกิด Overfitting นอกจากนี้ยังใช้เทคนิค Transfer Learning ร่วมกับการแช่แข็งเลเยอร์บางส่วน (Partial layer freezing) เพื่อเพิ่มประสิทธิภาพการฝึกสอนโมเดล

ผลการทดลองแสดงให้เห็นว่าวิธีการแบบสองขั้นตอนมีประสิทธิภาพดีกว่าวิธีการแบบขั้นตอนเดียว (One-stage) และโมเดลพื้นฐานอื่นๆ โดยมีความทนทานต่อภาพเบลอและมุมกล้องที่เอียง ระบบนี้จึงเป็นเครื่องมือที่มีศักยภาพในการคุ้มครองผู้บริโภคและตรวจสอบความปลอดภัยของอาหาร

\section{1. บทนำ (Introduction)}
ปัญหาการปลอมแปลงอาหารทะเล โดยเฉพาะการนำปลาเรนโบว์เทราต์มาขายในชื่อปลาแซลมอน เป็นปัญหาที่แพร่หลาย ซึ่งปลาเทราต์มีความเสี่ยงต่อพยาธิสูงกว่าและไม่เหมาะสำหรับการกินดิบ (เช่น ซาชิมิ) ความคล้ายคลึงทางกายภาพของเนื้อปลาทั้งสองชนิดทำให้ผู้บริโภคแยกแยะได้ยาก

\subsection{ความแตกต่างทางสายตา (Visual Similarities)}
\begin{itemize}
    \item \textbf{ปลาแซลมอน}: สีส้มแดงสดถึงชมพูเข้ม มีลายไขมันสีขาวแทรกชัดเจน (Marbling) เนื้อมีความมันวาว
    \item \textbf{ปลาเรนโบว์เทราต์}: สีส้มอ่อนถึงชมพูซีด ลายไขมันไม่ชัดเจน เนื้อมีความมันวาวน้อยกว่า
\end{itemize}

\subsection{วัตถุประสงค์ของงานวิจัย}
เพื่อพัฒนาระบบที่ใช้งานบนสมาร์ทโฟน ให้ผู้บริโภคสามารถถ่ายภาพเนื้อปลาและตรวจสอบได้ทันทีว่าเป็นปลาแซลมอนจริงหรือไม่ และมีความสดอยู่ในระดับใด

\section{2. ระเบียบวิธีวิจัย (Materials and Methods)}

\subsection{ชุดข้อมูล (Dataset)}
ภาพเนื้อปลาถูกรวบรวมและแบ่งเป็น 6 ประเภทหลัก:
\begin{enumerate}
    \item แซลมอนป่า (Wild Salmon) - แบบสเต๊ก (Steak) และแบบแล่ (Fillet)
    \item แซลมอนเลี้ยง (Farmed Salmon) - แบบสเต๊กและแบบแล่
    \item ปลาเทราต์ (Trout Meat)
    \item เนื้อสัตว์อื่นๆ (Other Meat)
\end{enumerate}
สำหรับการเกรดความสด แบ่งเป็น 3 ระดับตามสี: Pink Orange, Bright Orange (ดีที่สุด), Red Orange

\subsection{โมเดล Deep Learning}
งานวิจัยเลือกใช้ \textbf{DenseNet121} เป็นโครงสร้างหลัก เนื่องจากมีประสิทธิภาพในการเรียนรู้ฟีเจอร์ที่ซับซ้อนและใช้พารามิเตอร์น้อยกว่าโมเดลอื่น
\begin{itemize}
    \item \textbf{Improved DenseNet121}: มีการปรับปรุงโดยเพิ่ม Global Average Pooling (GAP) เพื่อลดมิติข้อมูล และเพิ่ม Dropout layer (0.2) เพื่อป้องกัน Overfitting
    \item \textbf{Transfer Learning}: ใช้เทคนิคการถ่ายโอนการเรียนรู้ 2 ขั้นตอน (ImageNet $\rightarrow$ Fish Classification $\rightarrow$ Freshness Grading) และมีการทดลองแช่แข็งเลเยอร์ (Freezing) ที่ 40\% เพื่อความสมดุลระหว่างความแม่นยำและเวลาในการเทรน
\end{itemize}

\subsection{กระบวนการทำงาน (System Workflow)}
\begin{enumerate}
    \item \textbf{Stage 1: Classification}: จำแนกภาพถ่ายออกเป็น 6 ประเภท (แซลมอนป่า/เลี้ยง สเต๊ก/แล่, เทราต์, อื่นๆ)
    \item \textbf{Stage 2: Freshness Grading}: หากระบุว่าเป็นแซลมอน จะทำการเกรดความสดต่อ (3 ระดับ)
\end{enumerate}

\section{3. ผลการทดลอง (Results)}

\subsection{ประสิทธิภาพการจำแนกประเภท (Classification Performance)}
โมเดล \textbf{Improved DenseNet121} ให้ผลลัพธ์ดีที่สุดเมื่อเทียบกับโมเดลอื่นๆ (เช่น ResNet, Inception, MobileNet):
\begin{itemize}
    \item \textbf{Accuracy}: 84.58\% (สูงที่สุด)
    \item \textbf{F1-Score}: สูงกว่า 77\% สำหรับ Wild Salmon Steak
\end{itemize}
การใช้ Transfer Learning โดยแช่แข็งเลเยอร์ 40\% ช่วยลดเวลาการเทรนลงได้ประมาณ 31\% โดยที่ความแม่นยำลดลงเพียงเล็กน้อย (ไม่มีนัยสำคัญทางสถิติ)

\subsection{ประสิทธิภาพการเกรดความสด (Freshness Grading Performance)}
สำหรับการเกรดความสด โมเดล Improved DenseNet121 ที่ใช้ Transfer Learning จากโมเดลจำแนกปลา (0\% freezing) ให้ความแม่นยำสูงสุดที่ \textbf{73.75\%} ซึ่งสูงกว่าการใช้ Machine Learning แบบดั้งเดิม (BPN, Fuzzy, ANFIS)

\subsection{การเปรียบเทียบ One-stage vs Two-stage}
\begin{itemize}
    \item \textbf{Two-stage} (แยกชนิดก่อน แล้วค่อยเกรดความสด): แม่นยำกว่า
    \item \textbf{One-stage} (ทำพร้อมกัน 14 คลาส): ความแม่นยำลดลงเหลือ 66.50\% เนื่องจากความซับซ้อนของคลาสที่มากขึ้น
\end{itemize}

\subsection{ความทนทาน (Robustness Analysis)}
\begin{itemize}
    \item \textbf{ภาพเบลอ (Motion Blur)}: โมเดลยังคงความแม่นยำได้สูงกว่า 80\% แม้ภาพจะเบลอเล็กน้อยถึงปานกลาง
    \item \textbf{มุมกล้องเอียง (Camera Tilt)}: ทนทานต่อการเอียงในแนวหน้า-หลัง ได้ดี แต่การเอียงซ้าย-ขวามากเกินไป (เกิน 8 องศา) ส่งผลกระทบต่อความแม่นยำอย่างมีนัยสำคัญ
\end{itemize}

\section{4. สรุปผล (Conclusion)}
งานวิจัยนี้นำเสนอระบบที่มีประสิทธิภาพและใช้งานได้จริงสำหรับผู้บริโภคในการตรวจสอบการปลอมแปลงปลาแซลมอนและตรวจสอบความสด ผ่านสมาร์ทโฟน
\begin{itemize}
    \item โมเดล Improved DenseNet121 แบบ Two-stage ให้ความแม่นยำสูงและมีประสิทธิภาพการคำนวณที่ดี
    \item ระบบมีความทนทานต่อสภาพแวดล้อมการถ่ายภาพจริงในระดับหนึ่ง
    \item ข้อจำกัดคือชุดข้อมูลยังค่อนข้างจำกัดและพึ่งพาการระบุป้ายกำกับ (Label) ด้วยสายตาผู้เชี่ยวชาญ งานวิจัยในอนาคตควรเพิ่มการตรวจสอบด้วยวิธีทางชีวเคมี
\end{itemize}

\end{document}
